\documentclass[12pt,a4paper]{letter}

\usepackage[T1]{fontenc}
\usepackage[utf8]{inputenc}
\usepackage{geometry}
\usepackage{url}
\usepackage[expansion=false]{microtype}

\geometry{
  top=2.5cm, bottom=2.5cm,
  left=2.8cm, right=2.8cm
}

\pagestyle{empty}

\address{
  Florin Olariu and Lenu\c{t}a Alboaie \\
  Department of Computer Science \\
  Alexandru Ioan Cuza University \\
  Ia\c{s}i, Romania \\
  \texttt{florin.olariu@info.uaic.ro; lenuta.alboaie@info.uaic.ro}
}

\date{September 2026}

\begin{document}

\begin{letter}{
  Prof. Song Guo, Editor-in-Chief \\
  \textit{IEEE Transactions on Cloud Computing} \\
  IEEE Computer Society
}

\opening{Dear Prof. Guo,}

Please find enclosed our revised manuscript, TCC-2026-04-0292,

\begin{center}
  \textbf{Transparent Multi-Cloud Migration Planning:\\
  A CSP--Expert--Pareto Decision Pipeline}
\end{center}

\noindent
by Florin Olariu and Lenu\c{t}a Alboaie, resubmitted in response to the
Major Revision decision. We thank the two reviewers for their careful and
constructive comments, which have materially improved the paper. A detailed,
point-by-point response to every comment is provided in the accompanying
``Response to Reviewers'' document; below we summarize the main changes.

\bigskip
\textbf{Summary of the revision.}
\begin{itemize}
  \item \textbf{No overstated optimizer claims.} Following Reviewer 1
        (Comment 1) and Reviewer 2, we no longer imply that CMO is a
        stronger optimizer than the baselines on individual metrics. We
        report plainly that a tuned Genetic Algorithm can match or exceed
        CMO on raw cost and latency, and reposition CMO's contribution
        around interpretability, auditability, and multi-criteria decision
        support rather than raw optimization performance.
  \item \textbf{Fair, like-for-like comparison.} Following Reviewer 1
        (Comment 2) and Reviewer 2, all baselines are now also scored with
        the same six-dimensional MCDA scheme used to rank CMO's solutions
        (cost, latency, reliability, security, vendor risk, scalability),
        so every comparison in the paper states explicitly which objective
        is being compared.
  \item \textbf{Honest scope on interpretability claims.} Following
        Reviewer 1 (Comment 3), we no longer imply that auditability and
        vendor-lock-in mitigation have been validated with users. We added
        a concrete, verbatim worked example of the constraint-proof
        artefact (Section V) and, drawing on Doshi-Velez and Kim's
        established taxonomy of interpretability evaluation, we explain in
        the Response to Reviewers why the present functionally-grounded
        evidence is the appropriate bar for this revision and why a small,
        non-expert pilot would not properly answer the application-grounded
        question being asked; a properly resourced practitioner study is
        scoped as concrete future work. The manuscript itself makes no
        unsupported claim in this regard.
  \item \textbf{Baseline configuration details.} Following Reviewer 1
        (Comment 4), we removed an unexecuted, placeholder NSGA-II/MOEA-D
        comparison that had been included in error, and we now document
        the Genetic Algorithm's encoding, population size, and operator
        settings in full so the remaining evolutionary baseline is
        reproducible.
  \item \textbf{Streamlined presentation.} Following Reviewer 1
        (Comment 5), the two internal configurations formerly labelled
        CMOv3 and CMOv4 are now presented explicitly as two configurations
        of a single pipeline (differing only in how latency is modelled),
        not as sequential software releases, to keep the focus on the
        research contribution rather than on version history.
  \item \textbf{Condensed abstract.} Following Reviewer 1 (Comment 6), the
        abstract has been shortened by nearly half and no longer lists
        specific run counts, timings, or component counts, focusing
        instead on motivation, design, and the main qualitative findings.
  \item \textbf{Clarified feasibility vs.\ completeness.} Following
        Reviewer 2, we now state explicitly that the constraint layer
        guarantees the feasibility of every returned solution, and we
        clarify the conditions under which the reported Pareto set is, or
        is not, claimed to be complete.
  \item \textbf{Weight sensitivity.} Following Reviewer 2, we added a new
        analysis (Section IV-D) showing how the final MCDA ranking responds
        to changes in the stakeholder weight profile, together with the
        rationale for the default weights used elsewhere in the paper.
  \item \textbf{Modest framing of the Netflix case study.} Following
        Reviewer 2, the Netflix-inspired scenario is now presented as an
        illustrative, publicly reproducible example grounded in public
        documentation and pricing, not as validation against real
        production constraints.
\end{itemize}

\bigskip
\textbf{Reproducibility.}
The implementation, experiment scripts, and result artefacts referenced in
the manuscript are maintained in a repository that is now publicly
available, as cited in the bibliography.

\bigskip
\textbf{Declarations.}
We confirm that the revised manuscript is original, has not been published
previously, and is not under consideration by any other journal or
conference. All authors have approved the submitted version, and there are
no conflicts of interest to declare.

\bigskip
We thank you and the reviewers again for the thorough evaluation, and we
believe the revised manuscript addresses all of the points raised. Please
let us know if any further information or clarification is required.

\closing{Yours sincerely,}

\vspace{0.5cm}
Florin Olariu \quad (Corresponding Author) \\
Lenu\c{t}a Alboaie \quad (Corresponding Author) \\[4pt]
Department of Computer Science, Alexandru Ioan Cuza University, Ia\c{s}i, Romania \\[2pt]
\texttt{florin.olariu@info.uaic.ro; lenuta.alboaie@info.uaic.ro}

\end{letter}
\end{document}
